% !TEX program = xelatex
% 공통수학2 · Ⅱ. 집합과 명제 · 1. 집합 · 02 집합 사이의 포함 관계 · 1/1차시
% 학생용 활동지 (답안 없음)
\documentclass[11pt,a4paper]{article}
\usepackage{kotex}
\usepackage{iftex}
\ifPDFTeX\else
  \setmainhangulfont{Noto Serif CJK KR}
  \setsanshangulfont{Noto Sans CJK KR}
\fi
\usepackage[margin=2.1cm]{geometry}
\usepackage{parskip}
\usepackage{enumitem}
\usepackage{array}
\usepackage{tabularx}
\usepackage{xcolor}
\usepackage{amssymb}
\usepackage{tikz}
\usepackage[most]{tcolorbox}
\usepackage{fancyhdr}
\usepackage{titlesec}

\definecolor{goldc}{HTML}{B07C22}
\definecolor{goldstrong}{HTML}{8A5F16}
\definecolor{indigoc}{HTML}{2B3B57}
\definecolor{indigosoft}{HTML}{6E82A6}
\definecolor{linec}{HTML}{D6DACB}
\definecolor{surface2}{HTML}{E4E8DC}
\definecolor{writeline}{HTML}{C7CCBA}
\definecolor{inksoft}{HTML}{52604F}

\titleformat{\section}{\normalfont\sffamily\bfseries\large\color{indigoc}}{\thesection}{0.6em}{}
\titlespacing{\section}{0pt}{14pt}{6pt}
\newtcolorbox{goalbox}{enhanced, breakable, colback=white, colframe=goldc, boxrule=0.5pt, leftrule=3pt, arc=2pt, left=10pt, right=10pt, top=8pt, bottom=8pt}
\newtcolorbox{sheetbox}{enhanced, breakable, colback=white, colframe=linec, boxrule=0.6pt, arc=3pt, left=11pt, right=11pt, top=9pt, bottom=9pt}
\newcommand{\writespace}[1]{%
  \par\nobreak\vspace{3pt}
  \foreach \n in {1,...,#1} {%
    \noindent\textcolor{writeline}{\rule{\linewidth}{0.4pt}}\par\vspace{0.62cm}%
  }
}
\newcommand{\fillblank}[1]{\makebox[#1]{\hrulefill}}
\pagestyle{fancy}\fancyhf{}\renewcommand{\headrulewidth}{0pt}
\fancyfoot[L]{\footnotesize\color{inksoft} 공통수학2 · 2026학년도 2학기 · 지평선고등학교}
\fancyfoot[R]{\footnotesize\color{inksoft} 다음 시간 --- 03 교집합과 합집합(1)}

\begin{document}

\noindent
\begin{minipage}[b]{0.62\linewidth}
  {\sffamily\footnotesize\color{indigosoft} 공통수학2 · Ⅱ. 집합과 명제 · 1. 집합}\\[2pt]
  {\sffamily\bfseries\Large 02 집합 사이의 포함 관계 \textperiodcentered\ 20차시}
\end{minipage}\hfill
\begin{minipage}[b]{0.36\linewidth}
  \raggedleft\small\color{inksoft}
  반\ \fillblank{1.6cm} \quad 번호\ \fillblank{1.6cm}\\[4pt]
  이름\ \fillblank{3.6cm}
\end{minipage}

\vspace{4pt}
\noindent{\color{goldc}\rule{\linewidth}{2.2pt}}
\vspace{10pt}

\begin{goalbox}
\textbf{오늘의 목표}\\
부분집합, 진부분집합, 서로 같은 집합의 뜻을 알고 두 집합 사이의 포함 관계를 판단할 수 있다.
\vspace{4pt}
\small\color{inksoft}
$\square$ 자신 있음 \quad $\square$ 조금 헷갈림 \quad $\square$ 처음 봄 \hfill (수업 전 나의 상태에 표시)
\end{goalbox}

\section{생각 열기 --- 악기 두 집합}

\begin{sheetbox}
피아노 삼중주 악기 $A=\{$피아노, 바이올린, 첼로$\}$ \\
피아노 사중주 악기 $B=\{$피아노, 바이올린, 비올라, 첼로$\}$

\vspace{4pt}
집합 $A$의 원소가 모두 집합 $B$에 속하는지 확인해 보자.
\writespace{2}
\end{sheetbox}

\section{개념 정리 --- 부분집합·서로 같음·진부분집합}

\begin{sheetbox}
두 집합 $A$, $B$에 대하여 $A$의 모든 원소가 $B$에 속할 때, $A$를 $B$의 \textbf{부분집합}이라 하고 $A\ \fillblank{1cm}\ B$로 나타낸다.

\vspace{6pt}
두 집합 $A$, $B$에 대하여 $A\subset B$이고 $B\subset A$일 때, `$A$와 $B$는 서로 같다'고 하며 $A\ \fillblank{1cm}\ B$로 나타낸다.

\vspace{6pt}
$A\subset B$이고 $A\neq B$일 때, $A$를 $B$의 \fillblank{3.2cm}이라고 한다.

\vspace{8pt}
{\footnotesize\color{inksoft} 주의: 모든 집합은 자기 자신의 부분집합이다. 또, 공집합은 모든 집합의 부분집합이다.}
\end{sheetbox}

\section{연습 문제}

\begin{sheetbox}
\textbf{1.} 다음 $\square$ 안에 $\subset$, $\not\subset$ 중 알맞은 것을 써넣으시오.\\
\hspace*{1.6em}⑴ $\{2,3,5\}\ \square\ \{x\mid x$는 10 이하의 소수$\}$\\
\hspace*{1.6em}⑵ $\{0,1\}\ \square\ \{x\mid x^2+x=0\}$
\writespace{3}

\vspace{6pt}
\textbf{2.} 다음 $\square$ 안에 $=$, $\neq$ 중 알맞은 것을 써넣으시오.\\
\hspace*{1.6em}⑴ $\{2,4,6\}\ \square\ \{x\mid x$는 8 이하의 짝수$\}$\\
\hspace*{1.6em}⑵ $\{x\mid x^2-1=0\}\ \square\ \{-1,1\}$
\writespace{3}

\vspace{6pt}
\textbf{3.} 집합 $A=\{a,b,c\}$에 대하여 다음에 답하시오.\\
\hspace*{1.6em}⑴ 집합 $A$의 부분집합을 모두 구하고, 그 개수를 말하시오.\\
\hspace*{1.6em}⑵ 집합 $A$의 진부분집합을 모두 구하고, 그 개수를 말하시오.
\writespace{5}

\vspace{6pt}
\textbf{4.} 평행사변형 전체의 집합 $A$, 직사각형 전체의 집합 $B$, 마름모 전체의 집합 $C$, 정사각형 전체의 집합 $D$가 있다. 이 네 집합 사이의 포함 관계를 모두 말해 보자.
\writespace{3}
\end{sheetbox}

\section{사고의 가시화 --- See · Think · Wonder}

\noindent{\footnotesize\color{inksoft} 오늘 배운 `집합 사이의 포함 관계'를 떠올리며 아래 세 칸을 채워 보자.}\\[6pt]
\begin{tabularx}{\linewidth}{@{}XXX@{}}
\textbf{\footnotesize\color{indigoc} See --- 무엇을 보았나?} &
\textbf{\footnotesize\color{indigoc} Think --- 무엇을 생각했나?} &
\textbf{\footnotesize\color{indigoc} Wonder --- 무엇이 궁금한가?} \\[4pt]
\end{tabularx}
\noindent
\begin{minipage}[t]{0.32\linewidth}\writespace{3}\end{minipage}\hfill
\begin{minipage}[t]{0.32\linewidth}\writespace{3}\end{minipage}\hfill
\begin{minipage}[t]{0.32\linewidth}\writespace{3}\end{minipage}

\section{마무리 성찰}

\begin{sheetbox}
\begin{minipage}[t]{0.48\linewidth}
{\footnotesize\color{inksoft} 오늘 배운 것을 한 문장으로}
\writespace{2}
\end{minipage}\hfill
\begin{minipage}[t]{0.48\linewidth}
{\footnotesize\color{inksoft} 감사 일기 / 유념 한 줄}
\writespace{2}
\end{minipage}
\end{sheetbox}

\end{document}
